% Original applied explanation; see Chapter04/SOURCES.md.
\subsubsection{তালিকার ধরন বাছাইয়ে একটি সহজ পরীক্ষা}
ভর্তি-আবেদনের ধাপ উল্টে দিলে কাজ বদলে যায়; সেখানে \texttt{ol} উপযোগী। প্রয়োজনীয় কাগজপত্রের তালিকায় ক্রম সাধারণত জরুরি নয়; সেখানে \texttt{ul} উপযোগী। পরিভাষা ও তার ব্যাখ্যা জোড়া আকারে রাখতে \texttt{dl}, \texttt{dt} ও \texttt{dd} লাগে।

বিভাগের তালিকার ভেতরে বিষয়ের উপতালিকা বসালে উপতালিকাটি সংশ্লিষ্ট \texttt{li}-এর ভেতরে রাখবে। এতে কোন বিষয় কোন বিভাগের অধীন, তা গঠনে প্রকাশ পায়। শুধু স্পেস দিয়ে লেখা সরিয়ে একই সম্পর্ক নির্ভরযোগ্যভাবে প্রকাশ করা যায় না।
